% TOP_PROBLEM: {"id":30007760,"problem_number":"LOCAL-30007760","title":"Unemployment-insurance design","category_id":21,"category":{"id":21,"name":"economics","display_name":"Economics","description":"Open economic theory statements, published research agendas, and proposed scoped specifications; question provenance and historical priority are qualified per record.","slug":"economics","order_index":21,"created_at":"2026-10-08T00:00:00Z"},"status":"research_question","external_url":null,"legacy_ids":[],"published":false,"statement":"In the explicitly specified setting: How should unemployment-benefit levels and duration adapt to recessions, household liquidity, job quality, and employer responses? The mathematical statement above fixes the target, assumptions and scope of this catalogue formulation.","economics":{"original_id":249,"field":"Taxation and social insurance","kind":"design","formalization_basis":"adapted_research_question","source_status":"research_agenda","priority_status":"earliest_found","priority_note":"Earliest explicit agenda verified in this audit, not a claim of first historical publication. The mathematical scope is authored for this catalogue.","earliest_verified_reference":{"authors":["Thomas Le Barbanchon","Johannes F. Schmieder","Andrea Weber"],"title":"Job Search, Unemployment Insurance, and Active Labor Market Policies","year":2024,"url":"https://www.rfberlin.com/wp-content/uploads/2024/10/24024.pdf","doi":"10.3386/w32720","locator":"Sections 3.6.2, 3.7–3.9; printed pp. 110–111, 117–130, especially pp. 128–130","evidence_summary":"Explicit future-work passages call for formalizing separation fiscal externalities, quantifying aggregate-demand effects and extending optimal UI to behavioral agents and partial-employment programs."},"current_reference":{"authors":["Thomas Le Barbanchon","Johannes F. Schmieder","Andrea Weber"],"title":"Job Search, Unemployment Insurance, and Active Labor Market Policies","year":2024,"url":"https://www.rfberlin.com/wp-content/uploads/2024/10/24024.pdf","doi":"10.3386/w32720","locator":"Sections 3.6.2, 3.7–3.9; printed pp. 110–111, 117–130, especially pp. 128–130","evidence_summary":"Explicit future-work passages call for formalizing separation fiscal externalities, quantifying aggregate-demand effects and extending optimal UI to behavioral agents and partial-employment programs."},"formalization_note":"Catalogue-authored scoped formalization. Population, instruments, welfare objective and identification restrictions must be instantiated before claiming a resolution; source authors are not credited with these added assumptions."},"statusQualification":"Published question or research agenda verified at the cited locator. The mathematical specification is adapted for this catalogue and includes additional assumptions; source publication does not establish first-ever priority or certify this exact model remains unsolved. Catalogue-authored scoped formalization. Population, instruments, welfare objective and identification restrictions must be instantiated before claiming a resolution; source authors are not credited with these added assumptions.","statusReviewedAt":"2026-10-08"}
% FIRST_POSTED: 2026-10-08
% ENTRY_KIND: statement_only
% !TeX program = lualatex
\documentclass[11pt]{article}
\usepackage{fontspec}
\usepackage[a4paper,margin=25mm]{geometry}
\usepackage{amsmath,amssymb,mathtools}
\usepackage{xurl}
\usepackage[hidelinks]{hyperref}
\newcommand{\sref}[2]{\hyperlink{source:#1:#2}{[S#2]}}
\setlength{\emergencystretch}{3em}
\begin{document}
% problemId:30007760
\section*{Unemployment-insurance design}\label{30007760}
\subsection{Definitions and mathematical statement}
Specify a dynamic labor market with household types \(\theta\sim F\) for a fixed population distribution \(F\), assets \(a\), unemployment duration \(d\), aggregate state \(z\), discount factor \(\beta\in(0,1)\), and preferences \(u_\theta(c)-v_\theta(s)\) over consumption and search effort. Firms choose vacancies and separations; matching determines employment and wages. A benefit rule \(b(d,a,z,\theta_o)\) can depend only on duration, legally observable assets and characteristics \(\theta_o\). Let \(\mathcal B\) be a specified admissible rule set and \(\mathcal E(b)\) its nonempty equilibria, including household choices, employer responses and balanced government accounts. For fixed nonnegative welfare weights \(\omega_\theta\), normalized by \(\int\omega_\theta dF(\theta)=1\), characterize
\[\sup_{b\in\mathcal B}\inf_{e\in\mathcal E(b)}\mathbb E_e\!\left[\sum_{t\geq0}\beta^t\omega_\theta\{u_\theta(c_t)-v_\theta(s_t)\}\right].\]
The expectation averages over \(F\) and equilibrium histories, with convergent discounted welfare. The target includes estimating the consumption-insurance gain, search and separation fiscal externalities, accepted-job quality and aggregate-demand effects required for policy comparisons. Specify whether preferences used for welfare coincide with behavior when beliefs or present bias are modeled. Derive empirically implementable sufficient statistics or sharp bounds under explicit restrictions, and evaluate benefit-level and duration reforms in recession and expansion states. Comparison of individual unemployment durations alone is insufficient when vacancy creation, separations and uninsured workers respond. The formal scope is an adapted research agenda, not an existence conjecture for every matching economy.

\par\textbf{Formulation and scope.} Catalogue-authored scoped formalization. Population, instruments, welfare objective and identification restrictions must be instantiated before claiming a resolution; source authors are not credited with these added assumptions.

\par\textbf{Open-question provenance.} An explicit question or research agenda was verified at the source locator. This does not by itself certify that every later version remains unresolved \sref{30007760}{1}.

\par\textbf{Historical priority.} Earliest explicit agenda verified in this audit, not a claim of first historical publication. The mathematical scope is authored for this catalogue.

\subsection{Short English statement}
In the explicitly specified setting: How should unemployment-benefit levels and duration adapt to recessions, household liquidity, job quality, and employer responses? The mathematical statement above fixes the target, assumptions and scope of this catalogue formulation.

\subsection{Sources}
\begin{itemize}\raggedright
\item[\textnormal{[S1]}]\hypertarget{source:30007760:1}{} Thomas Le Barbanchon, Johannes F. Schmieder, Andrea Weber. 2024. Job Search, Unemployment Insurance, and Active Labor Market Policies. Sections 3.6.2, 3.7–3.9; printed pp. 110–111, 117–130, especially pp. 128–130.
\url{https://www.rfberlin.com/wp-content/uploads/2024/10/24024.pdf}.
DOI: 10.3386/w32720.
{\small\textit{Earliest verified question/agenda reference; historical priority qualified below. Explicit future-work passages call for formalizing separation fiscal externalities, quantifying aggregate-demand effects and extending optimal UI to behavioral agents and partial-employment programs.}}
\end{itemize}
\end{document}
